\documentclass[11pt]{article}
\usepackage[a4paper,margin=2.5cm]{geometry}
\usepackage[T1]{fontenc}
\usepackage{lmodern}
\usepackage{microtype}
\usepackage{amsmath}
\usepackage{booktabs}
\usepackage{tabularx}
\usepackage{array}
\usepackage[authoryear,round]{natbib}
\usepackage{xcolor}
\usepackage{listings}
\usepackage[hidelinks]{hyperref}

\lstset{basicstyle=\ttfamily\footnotesize,breaklines=true,frame=single,columns=fullflexible,
  keywordstyle=\bfseries,language=Python,showstringspaces=false}
\newcolumntype{L}{>{\raggedright\arraybackslash}X}

\title{Backtests That Do Not Survive: An Out-of-Sample Test of 41 Swing Trading Signals, Five Original Strategies and Funding-Rate Carry in U.S. Stocks, ETFs and Crypto}

\author{Telmen Nyamdelger\\[2pt]
\small Independent researcher\\[6pt]
\small with Claude (Anthropic), an AI model used as research engineer and reviewer\thanks{%
The study was designed and run jointly by the human author and Claude. Claude wrote and reviewed
the software and drafted the text under the author's direction; the author chose the research
questions, approved every protocol before results were seen, and takes responsibility for the
paper. See the contribution statement at the end.}}

\date{November 2026}

\begin{document}
\maketitle

\begin{abstract}
Retail traders are sold swing trading systems on the strength of attractive backtests. We built an
open-source research pipeline that tries to give such systems a fair chance and an honest test.
It covers 41 textbook technical signals with fixed parameters, pairs and triples of those signals,
a walk-forward gradient boosting model, a formal validation gate, a capacity-limited portfolio
backtest, five new strategies designed from patterns found in the data, a tactical trend rule and a
crypto funding-rate carry trade. The data are 888 U.S. large and mid-cap stocks, 100 ETFs and 419
Binance crypto assets including 63 delisted coins, with costs charged on every trade. Each test was
written down before its result was seen, and the most recent data (2024 to 2026) were locked away
and used once.
No strategy passed. Several looked excellent in sample: one crypto strategy returned 1,028\% over
2018 to 2021 and then lost 52\% over 2022 to 2023. The one signal that passed the 2018 to 2023
check, buying U.S. stocks when RSI(14) falls below 30, did worse than random stocks from the same
list and than the S\&P 500 ETF in 2025 to 2026. A trend rule cut drawdowns but earned less than
simply holding the same ETFs, and the carry trade was smooth but earned less than cash in 2022 to
2023. Simple buy-and-hold matched or beat every strategy we tested. We release all code and data so
that others can check these results and test their own ideas the same way.
\end{abstract}

\noindent\textbf{Keywords:} technical analysis, backtest overfitting, out-of-sample testing,
survivorship bias, multiple testing, cryptocurrency, retail trading

\noindent\textbf{Disclaimer:} This paper is research and education. It is not investment advice.

\section{Introduction}

Many people believe they can make a living from short-term trading rules: buy when a stock is
oversold, sell when a moving average crosses, follow a breakout. Courses, signal services and
trading bots are marketed with backtests that show smooth, rising equity curves. The evidence on
how retail traders actually do is far less flattering. \citet{barber2000trading} found that the
households who traded most earned the least. In Taiwan, fewer than 1\% of day traders earned
reliable positive returns after fees \citep{barber2014speculator}. In Brazil, 97\% of people who day
traded equity futures for more than 300 days lost money \citep{chague2019daytrading}.

Part of the gap between marketed backtests and real results is statistical. If enough rules or
parameter settings are tried, some will look excellent by chance, and their out-of-sample returns
can be negative \citep{bailey2014pseudo}. Finance research has the same problem at a larger scale:
hundreds of published return predictors have not survived stricter testing
\citep{harvey2016cross}. Survivorship bias adds to it. A stock universe built from today's index
members leaves out every company that failed, which makes buying weakness look better than it was
\citep{brown1992survivorship,elton1996survivor}.

This paper reports a long, honest attempt to find a swing trading strategy that works. The project
started as an attempt to build a profitable trading bot. It became a study of why backtests mislead
when no strategy survived. We wanted the strategies to have every fair chance, so the pipeline
covers single signals, combinations, machine learning, and strategies we designed ourselves from
patterns in the data. We also wanted the test to be hard to fool, so every protocol and pass
criterion was fixed in writing before its result was computed, and the newest data were held back
for a single final test.

Our contributions are three. First, we test 41 common swing signals across U.S. stocks, ETFs and
crypto with date-clustered, overlap-aware statistics, multiple-testing control and costs on every
trade, and we show how survivorship and regime dependence create false winners. Second, we take
the most promising results through a pre-registered validation gate, a portfolio backtest and a
one-time holdout, and show that none of them holds up. Third, we release the complete pipeline,
data snapshots and logs, so that anyone can repeat the work or test a new idea with the same
discipline. The project was carried out by a human researcher working with an AI model, and we
describe how that collaboration was organised and checked.

\section{Related literature}

Early evidence that simple technical rules predicted stock index returns \citep{brock1992simple} did
not survive a correction for data snooping across the universe of rules
\citep{sullivan1999data,white2000reality}. \citet{lo2000foundations} formalised chart patterns and
found some information in them, while the survey by \citet{park2007profitability} concluded that
apparent profits shrank over time and were sensitive to data snooping and costs.
Momentum is better documented: cross-sectional momentum in stocks \citep{jegadeesh1993returns} and
time-series momentum across asset classes \citep{moskowitz2012time}, although momentum strategies
suffer rare, severe crashes \citep{daniel2016momentum}. \citet{faber2007quantitative} proposed a
10-month moving average rule for tactical asset allocation that mainly reduces drawdowns. In crypto,
\citet{liu2021risks} find strong time-series momentum in coin returns, and \citet{he2022perpetual}
show that perpetual futures trade away from spot and that the gap, which funding payments are meant
to close, has narrowed over time.

On method, we follow the literature on backtest overfitting
\citep{bailey2014pseudo,bailey2014deflated,lopezdeprado2018advances}: limit the number of trials,
fix parameters in advance, and keep untouched data for a final test. We use Newey--West standard
errors for overlapping returns \citep{newey1987simple}, the Benjamini--Hochberg false discovery rate
\citep{benjamini1995controlling}, gradient boosting \citep{ke2017lightgbm} and SHAP attributions
\citep{lundberg2017unified}.

\section{Data}

\paragraph{U.S. stocks and ETFs.} Daily prices from Yahoo Finance, adjusted for splits and dividends,
from January 2000. The stock universe is the current S\&P 500 and S\&P 400 membership, 888 single
stocks after removing ETFs. Because it contains only today's survivors, we treat every stock result
as suspect and label it so in all reports. A separate universe of 100 liquid ETFs (U.S. equity,
factors, sectors and industries, international, bonds, commodities, currencies and real estate) has
no constituent survivorship and serves as the primary equity result.

\paragraph{Crypto.} Daily Binance spot data for every USDT pair ever listed (632 after removing
stablecoins, fiat, wrapped, tokenized-equity and leveraged tokens), taken from the public archive
so that delisted coins are included. Binance reuses tickers: \texttt{LUNAUSDT} was Terra LUNA
until it collapsed in May 2022 and a different coin afterwards, and both sources return one joined
series. We split any series with a gap longer than five days into separate assets, which affected
seven tickers. Of the 419 assets with data before 2024, 63 were delisted; a delisted coin's open
trades exit at its last close. An asset-day counts only if the coin ranks in the top 100 by
trailing 30-day median dollar volume at that day's close, a point-in-time liquidity filter.

\paragraph{Perpetual futures.} Binance USD-margined funding rates and daily perpetual prices for 469
contracts from January 2020, including delisted contracts.

\paragraph{Splits.} Table~\ref{tab:splits} shows the periods. Data from 1 January 2024 onward could
be loaded only through a logged flag and only by a final-test script that refuses to run twice.
Japanese (Nikkei 225) and European (DAX, CAC 40, FTSE 100) stocks were downloaded as a fresh test
market but never inspected, because no stock strategy reached the stage that would use them.

\begin{table}[h]
\centering\small
\caption{Sample periods.}\label{tab:splits}
\begin{tabular}{lll}
\toprule
 & Stocks and ETFs & Crypto \\
\midrule
Discovery (development) & 2000 to 2017 & 2018 to 2021 (carry: 2020 to 2021) \\
Validation & 2018 to 2023 & 2022 to 2023 \\
Holdout (used once) & 2024 to Nov.\ 2026 & 2024 to Nov.\ 2026 \\
\bottomrule
\end{tabular}
\end{table}

\section{Method}

\subsection{Signals and outcomes}
Each signal is a pure function of prices up to the close of day $t$ and returns true or false.
The 41 signals cover trend, pullback, breakout, mean reversion, momentum, divergence, volume,
support and resistance, and candlestick families, all with textbook parameters (for example
RSI(14) computed with Wilder smoothing, \citealp{wilder1978new}); Appendix~\ref{app:signals} lists
them. A truncation test confirms that no signal changes when future data are removed. Trades enter
at the open of day $t+1$. The main horizon is 10 trading days. Round-trip costs are 0.10\% for U.S.
stocks and ETFs and 0.30\% for crypto. Excess return is the trade's return net of costs minus the
benchmark's return over the same dates (SPY for stocks and ETFs, BTC for crypto). Returns are
winsorised at the 0.1st and 99.9th percentiles of each split before any statistic.

\subsection{Statistics}
Signals fire on many assets on the same day, and 10-day returns overlap. We therefore average per
date first and measure a signal's \emph{edge} as its same-day excess return minus the same-day
average of all assets in the same universe. The mean edge is tested with a Newey--West regression
on the date-sorted series with $2h$ lags and a $t$ distribution with conservative effective degrees
of freedom, $\max(\lfloor n_{\text{dates}}/2h\rfloor-1,1)$. Simulations with dense and sparse firing
patterns gave rejection rates of about 4 to 8\% at the nominal 5\% level. Benjamini--Hochberg
$q$-values control the false discovery rate within each table. Rows with fewer than 100 occurrences
or $\max(100,5h)$ distinct dates are marked insufficient and never ranked.

A signal is a \emph{candidate} only if it is sufficient, $q<0.10$, its edge is positive, its mean
excess return over the benchmark is positive, and its mean absolute return net of costs is
positive. The last condition was added after a review found that short signals in crypto could beat
BTC while losing money as plain shorts.

\subsection{Combinations and machine learning}
For every pair of same-direction signals we measure \emph{lift}: the edge of days when both fire
minus the better of the two single-only edges, tested with a conservative two-sample HAC test.
Triples are formed only from pairs with positive lift. A redundancy map flags signals that fire
on the same rows. For machine learning, a LightGBM regressor with fixed hyperparameters predicts
10-day excess returns from the signal matrix and regime labels. It is trained walk-forward by
calendar year with an embargo of $h+1$ trading days, scored on the validation split, and explained
with SHAP values.

\subsection{Validation gate, portfolio test and holdout}
Discovery candidates and synergistic pairs are re-tested on the validation split with the same
rule. The portfolio test holds at most 10 positions of equal size for 10 trading days, enters at
the next open, fills free slots in random order (five seeds) and books each trade's return net of
costs. A same-universe random baseline (enter any eligible row with the same slots and seeds)
measures what the universe alone would have earned, which is the fair comparison when the universe
itself is biased.

Every later experiment (the RSI test on 2025 to 2026, the trend rule, the five original strategies
and the carry trade) was specified in a dated, committed document before results were computed,
including the pass criteria. A holdout run writes a record file before loading any 2024 or later
data, logs every access, and refuses to run a second time.

\subsection{How the human and AI collaboration was organised}
The human author set the goals and made every decision about scope, markets and when to use the
holdout. Claude wrote the code in small tasks. Each task was implemented by one AI agent and
checked by a separate, fresh AI reviewer against the written specification before the next task
began, and each phase ended with a whole-code review. These reviews found real errors that would
have changed the conclusions, among them a price filter that used future information (Yahoo's
back-adjusted prices are lower in the past for stocks that later rose, so dropping sub-\$1 rows
removed early Nvidia and Apple data), the LUNA ticker splice, a funding-payment timing error that
favoured one carry strategy, and the short-signal issue above.

\section{Results}

\subsection{Single signals}
Table~\ref{tab:solo} summarises the discovery results at the 10-day horizon. No signal had a
significant edge in the 100-ETF universe. In crypto, trend and breakout signals were candidates and
every dip-buying signal had a significantly negative edge. In U.S. stocks the pattern reversed: seven
dip-buying signals were candidates and seven breakout and trend signals were significantly worse
than average.

\begin{table}[h]
\centering\small
\caption{Discovery results for 41 signals, 10-day horizon. Edge is in basis points per trade over the
same-day universe average.}\label{tab:solo}
\begin{tabularx}{\textwidth}{lccL}
\toprule
Universe & Candidates & Significantly negative & Strongest results (edge, $q$) \\
\midrule
ETFs (100) & 0 & 0 & none \\
Crypto (419) & 5 & 6 & MA 20/50 cross up (+303, 0.056); Donchian 55 breakout (+286, 0.056);
lower Bollinger close ($-$347, 0.0002) \\
U.S. stocks (888, suspect) & 7 & 7 & RSI(14) oversold (+44, 0.002); stochastic oversold cross (+35,
0.0007); Donchian 20 breakdown ($-$44, $<$0.001) \\
\bottomrule
\end{tabularx}
\end{table}

The stock candidates should be read against the universe's drift. Across all discovery rows, the
stock universe beat SPY by 33.9 basis points per 10 days on average, which is mostly survivorship.
The ETF universe, which does not have this problem, showed nothing.

\subsection{Combinations and machine learning}
Of 582 same-direction pairs per universe, none showed significant synergy in ETFs or crypto. In
U.S. stocks two pairs showed synergy and five showed conflict, all with $q=0.0998$, just under the
threshold. One of the two synergistic pairs (MACD bullish divergence with OBV leading price)
later had a significantly \emph{negative} edge on validation ($-$91 basis points, $q=0.001$).

Table~\ref{tab:ml} shows the walk-forward model. In ETFs it had no predictive power. In crypto only
one walk-forward year could be scored, and its validation spread was zero. In stocks it had a small
positive spread, and SHAP attributed most of it to down-trend and oversold features, the same
pattern that survivorship inflates.

\begin{table}[h]
\centering\small
\caption{Walk-forward LightGBM. IC is the mean daily rank correlation between prediction and
outcome; spread is the top-decile excess return over the daily average, in basis points.}\label{tab:ml}
\begin{tabular}{lrrrr}
\toprule
 & IC (walk-forward) & Spread (walk-forward) & IC (validation) & Spread (validation) \\
\midrule
ETFs & 0.007 ($p=0.42$) & $-$0.5 & 0.029 ($p=0.047$) & $-$1.5 \\
Crypto & 0.029 ($p=0.053$) & 119 & 0.024 ($p=0.060$) & $-$2.4 \\
U.S. stocks & 0.010 ($p=0.006$) & 11 & 0.010 ($p=0.12$) & 22 \\
\bottomrule
\end{tabular}
\end{table}

\subsection{Validation gate and portfolio test}
None of the five crypto candidates passed validation: four had negative edges in 2022 to 2023. Of
the seven stock candidates, only RSI(14) oversold passed (edge +47 basis points, $q=0.041$). As a
portfolio over 2018 to 2023 it returned a median of 151\% across 15 runs (range 91\% to 251\%),
against 84\% for the same-universe random baseline and 96\% for SPY. Its maximum drawdown was 47\%
against 34\% for SPY, and its Sharpe ratio (0.72 in the example run) was no better than SPY's (0.71).
Because the stock list is biased, it was not eligible for the holdout under our protocol.

\subsection{Holdout and follow-up tests}

\paragraph{RSI(14) oversold, 2025 to 2026.} At the author's request, the stock holdout was used once
to test this signal unchanged on January 2025 to November 2026. It failed. Its edge over the same-day
average was $-$13.4 basis points ($p=0.70$). As a 10-slot portfolio without per-trade minimum fees
it returned a median of 6.5\%, against 11.2\% for random stocks from the same list and 24.9\% for
SPY. We also simulated a \$20 paper account with fractional shares, a one-cent regulatory fee per
sale and optional 2$\times$ leverage. Every configuration lost money (median end balance \$11 to
\$17), the cent fee alone cost 15 to 50 basis points per trade, and leverage made every case worse.
U.S. brokers also require \$2,000 for a margin account.

\paragraph{Trend rule on ETFs.} We tested Faber's 10-month rule on five ETFs (SPY, EFA, IEF, VNQ,
DBC, with SHY as cash), pass bar fixed in advance: positive return, Sharpe at least that of holding
the same ETFs, and a smaller maximum drawdown. Table~\ref{tab:trend} shows that it passed in both
earlier periods only because it cut drawdowns, and was not confirmed on the holdout, where it
earned 22.6\% against 34.3\% for holding the same ETFs.

\begin{table}[h]
\centering\small
\caption{Faber 10-month trend rule versus holding the same five ETFs.}\label{tab:trend}
\begin{tabular}{lrrrrrr}
\toprule
 & \multicolumn{3}{c}{Trend rule} & \multicolumn{3}{c}{Hold same ETFs} \\
\cmidrule(lr){2-4}\cmidrule(lr){5-7}
Period & Return & Max DD & Sharpe & Return & Max DD & Sharpe \\
\midrule
2006 to 2017 & 61\% & $-$14\% & 0.67 & 60\% & $-$47\% & 0.40 \\
2018 to 2023 & 28\% & $-$10\% & 0.64 & 43\% & $-$26\% & 0.52 \\
2024 to 2026 (holdout) & 23\% & $-$6\% & 1.07 & 34\% & $-$11\% & 1.33 \\
\bottomrule
\end{tabular}
\end{table}

\paragraph{Five original strategies.} We then designed five strategies from patterns seen in the
earlier results and fixed their rules in advance: crypto momentum rotation with a BTC trend switch
(C1), a crypto weakness hedge (C2), C1 gated by the Crypto Fear \& Greed Index (C3), dip buying
inside uptrends (S1) and monthly stock momentum with a crash switch (S2). The final tests were
reserved for crypto 2024 to 2026 and for the unused Japanese and European stocks.
Table~\ref{tab:orig} shows that all five failed validation, so neither final test ran.

\begin{table}[h]
\centering\small
\caption{Original strategies. Benchmarks: BTC buy-and-hold (C1, C3); absolute (C2); same-universe
random entries (S1); equal-weight holding of the same stocks (S2).}\label{tab:orig}
\begin{tabularx}{\textwidth}{lrrrL}
\toprule
 & \multicolumn{2}{c}{Total return} & & \\
\cmidrule(lr){2-3}
Strategy & Discovery & Validation & Val.\ Sharpe & Reason for failing validation \\
\midrule
C1 crypto momentum & 117\% & $-$8\% & 0.06 & Sharpe below BTC (0.19) \\
C2 crypto weakness hedge & $-$16\% & 19\% & 0.46 & Sharpe below 0.5 \\
C3 mood + momentum & 1,028\% & $-$52\% & $-$0.78 & Sharpe below BTC \\
S1 dip in uptrend & 523\% & 37\% & 0.44 & Did not beat random entries \\
S2 momentum + crash switch & 1,753\% & 135\% & 0.88 & Below equal-weight hold (196\%) \\
\bottomrule
\end{tabularx}
\end{table}

C3 is the clearest example of overfitting to one regime. Its only data-chosen setting, whether to
buy when the index showed fear or greed, was picked on 2018 to 2021, a period dominated by one
crypto boom. In 2022 to 2023 it bought the strongest coins during short bear-market rallies.

\paragraph{Funding-rate carry.} The last test was a trade with an economic reason rather than a
price forecast: hold a coin and short the same notional in its perpetual future, and collect the
funding that leveraged longs pay. We tested a static BTC and ETH version (F1) and a weekly rotation
into the five coins with the highest recent funding (F2). The pass bar was an annual return above
4\%, a Sharpe ratio of at least 1.0 and a maximum drawdown under 15\%.

\begin{table}[h]
\centering\small
\caption{Funding-rate carry, costs included.}\label{tab:carry}
\begin{tabular}{llrrrr}
\toprule
Strategy & Period & Annual return & Sharpe & Max DD & Turnover per year \\
\midrule
F1 BTC+ETH & 2020 to 2021 & 15.6\% & 3.66 & $-$0.5\% & 5 \\
F1 BTC+ETH & 2022 to 2023 & 2.5\% & 3.25 & $-$0.4\% & 4 \\
F2 rotating & 2020 to 2021 & 20.5\% & 2.76 & $-$0.6\% & 72 \\
F2 rotating & 2022 to 2023 & $-$5.5\% & $-$0.91 & $-$12.5\% & 76 \\
\bottomrule
\end{tabular}
\end{table}

F1 behaved as carry theory predicts, with almost no losses, but in the 2022 to 2023 bear market
demand for leverage fell and it earned 2.5\% a year, less than short-term Treasury bills paid at the
time. F2's turnover consumed its funding income. Both failed validation, so the crypto holdout is
still unused.

\section{Discussion}

The results repeat a familiar lesson with unusual clarity. Results that look strong in sample are
usually not real. C3 earned 1,028\% in development; the stock momentum strategy S2 earned 1,753\%.
Neither survived the next period. The one stock signal that survived validation failed when it was
finally tested on new data, and the universe it was tested in favoured it.

Four mechanisms produced the false winners. Multiple testing: with 41 signals, 582 pairs per
universe and several strategy families, some results clear any fixed bar by chance, which is why
$q$-values and a held-out test matter. Survivorship: the stock universe beat SPY by 34 basis points
per 10 days before any signal was applied, and every stock "edge" was a dip-buying signal, which
is what survivorship rewards. Regime dependence: crypto breakouts and C3 worked in one boom and
failed in the following bear market, in line with the crash risk of momentum
\citep{daniel2016momentum}. Costs and frictions: turnover destroyed F2, and small accounts face
minimum fees that exceed the edge of every signal we tested.

For retail investors the practical message is simple. A backtest that was not tested on data it had
never seen is close to worthless, and one built on today's index members is biased toward buying
dips. Even honest, well-built strategies in our study did no better than holding a diversified
index, which is what most research on active retail trading also finds
\citep{barber2000trading,barber2014speculator,chague2019daytrading}.

The AI collaboration made a study of this size possible for one person, but it did not make the
study correct by itself. Several conclusions changed because a separate reviewer checked each piece
of code against its written specification. We think the combination of written protocols,
independent review and a locked holdout is what readers should look for in any AI-assisted
financial research.

\section{Limitations}

The data are free and have known defects: the U.S. stock universe is survivorship-biased and Yahoo
prices are back-adjusted. We used one main horizon (10 days) and textbook parameters only, so a
tuned version of a signal could behave differently, although tuning would also raise the risk of
overfitting. Costs are simplified (no taxes, no market impact, fixed per-trade costs). Some
validation results were examined informally before the formal gate was written, and the stock and
ETF holdouts were each opened once for the tests reported here; reproducing those numbers re-reads
data that have already been seen. The crypto holdout and the Japanese and European data remain
unused and are available for one future test.

\section{Conclusion}

We gave 41 common swing trading signals, their combinations, a machine learning model, five new
strategies, a trend rule and a carry trade a fair and careful test on stocks, ETFs and crypto. None
passed. Several earned remarkable returns in the data they were built on and failed afterwards.
Simply holding the market matched or beat all of them. We hope the released pipeline helps traders
check an idea honestly before they risk money on it.

\section*{Data and code availability}
All source code, tests, protocols and the result files used in the tables are public at
\url{https://github.com/Temka921/tradingbot}. The \texttt{paper/data} folder holds the result
tables, and the \texttt{docs/superpowers/specs} folder holds each dated pre-registered design.
Prices can be re-downloaded with the included scripts.

\section*{Author contributions and use of AI}
Telmen Nyamdelger conceived the project, chose the research questions and markets, approved every
design and pass criterion before results were computed, decided when the holdout data could be used,
and is responsible for the paper. Claude (Anthropic) designed the statistical tests with the author,
wrote and reviewed the software, ran the experiments and drafted the manuscript, which the author
reviewed. Claude is not listed as an author because current publication policies do not allow AI
systems to take authorship responsibility; its role is disclosed here in full.

\bibliographystyle{plainnat}
\bibliography{references}

\appendix
\section{Signals tested}\label{app:signals}
\small
Short signals are marked with an asterisk; all other signals are long.
\begin{description}
\item[Trend (11):] MA 20/50 cross up, MA 20/50 cross down*, MA 50/200 cross up, MA 50/200 cross
down*, close crosses above MA 200, MACD cross up, MACD cross down*, ADX $>$ 25 with +DI above
$-$DI, Supertrend flip up, Supertrend flip down*, Parabolic SAR flip up.
\item[Pullback (4):] touch of the 20-day MA in an uptrend, RSI(14) between 40 and 50 in an uptrend,
RSI(2) below 10 above the 200-day MA, Fibonacci 50 to 61.8\% retracement.
\item[Breakout (6):] Donchian 20 breakout, Donchian 20 breakdown*, Donchian 55 breakout,
Bollinger/Keltner squeeze release, breakout on 1.5$\times$ volume, 52-week high.
\item[Mean reversion (7):] close below lower Bollinger band, close above upper band*, RSI(14) below
30, RSI(14) above 70*, z-score(20) below $-$2, stochastic oversold cross, gap-down fill setup.
\item[Momentum (2):] 12-1 month relative strength top decile, ROC(20) acceleration.
\item[Divergence (2):] RSI bullish divergence, MACD bullish divergence.
\item[Volume (3):] OBV new high before price, close above anchored VWAP from swing low,
accumulation/distribution rising.
\item[Support and resistance (2):] bounce off a prior swing low, pivot-point support hold.
\item[Candlesticks (4):] bullish engulfing, hammer, inside-bar breakout, morning star.
\end{description}

\section{Selected source code}
The full code is in the repository. The listings below are the core significance test and the
validation gate's decision rule.

\noindent\textbf{HAC test of a date-sorted series} (\texttt{src/swingbot/analysis/stats.py}):
\lstinputlisting[firstline=46,lastline=74]{../src/swingbot/analysis/stats.py}

\noindent\textbf{Same-day matched edge test} (\texttt{src/swingbot/analysis/stats.py}):
\lstinputlisting[firstline=101,lastline=113]{../src/swingbot/analysis/stats.py}

\noindent\textbf{Validation gate decision rule} (\texttt{src/swingbot/analysis/validation.py}):
\lstinputlisting[firstline=30,lastline=44]{../src/swingbot/analysis/validation.py}

\end{document}
